\documentclass[final]{beamer}

\usepackage[size=custom, width=106.68, height=106.68, scale=1.4]{beamerposter}
% 42in = 106.68cm

\usepackage{graphicx}
\usepackage{amsmath}
\usepackage{xcolor}
\usepackage{tikz}
\usepackage{tcolorbox}
\usetikzlibrary{calc}

%% ── GW brand colors ──────────────────────────────────────────
\definecolor{gwblue}{RGB}{0,59,92}
\definecolor{gwbuff}{RGB}{168,153,104}
\definecolor{gwcyan}{RGB}{0,151,215}
\definecolor{gwlightbg}{RGB}{237,229,210}   % gwbuff!35!white approx

%% ── Beamer theme: stripped down ──────────────────────────────
\usetheme{default}
\usecolortheme{default}
\setbeamercolor{background canvas}{bg=gwlightbg}
\setbeamertemplate{navigation symbols}{}

%% ── tcolorbox block style ────────────────────────────────────
\tcbuselibrary{skins}

\newenvironment{posterblock}[2][gwcyan]{%
  % #1 = header bg color, #2 = title
  \begin{tcolorbox}[
    enhanced,
    colback=white,
    colframe=#1,
    colbacktitle=#1,
    coltitle=white,
    fonttitle=\bfseries\large,
    title={#2},
    arc=4pt,
    boxrule=1.5pt,
    titlerule=0pt,
    top=6pt, bottom=6pt, left=8pt, right=8pt,
    toptitle=4pt, bottomtitle=4pt,
  ]
}{%
  \end{tcolorbox}
}

\DeclareMathOperator{\arctanh}{arctanh}

\begin{document}
\begin{frame}[t]{}

\begin{tikzpicture}[remember picture, overlay,
    every node/.style={inner sep=0, outer sep=0}]

  \coordinate (TL) at (current page.north west);

  %% ============================================================
  %% HEADER BAR
  %% ============================================================
  \fill[gwblue] (TL) rectangle ++(106.68cm, -14cm);

  % GWSPH Research Day logo — top left
  \node[anchor=north west]
    at ($(TL)+(1.5cm,-1.2cm)$) {
      
\includegraphics[height=8cm]{../figures/gwsph_research_day_logo.png}
    };

  % GW | Public Health logo — top right
  \node[anchor=north east]
    at ($(TL)+(105.5cm,-1.2cm)$) {
      
\includegraphics[height=10cm]{../figures/gw_public_health_logo.png}
    };

  % Title centered between logos
  \node[anchor=north, text width=68cm, align=center]
    at ($(TL)+(53.34cm,-0.6cm)$) {
      {\color{white}\fontsize{48}{56}\bfseries\selectfont
       \\ Computational Pipeline Choice Drives Shotgun-Amplicon\\[0.2ex]
       Discordance in Microbial Community Profiling:\\[0.2ex]
       A Fisher-Rao Distance Approach}\\[0.8ex]
      {\color{gwbuff!80!white}\fontsize{40}{42}\bfseries\selectfont
       Clark Gaylord \quad Marcos P\'{e}rez-Losada \quad Keith A.\ Crandall}\\[0.4ex]
      {\color{white!70!gwblue}\fontsize{36}{32}\selectfont
       Dept.\ of Biostatistics \& Bioinformatics and Computational Biology Institute \\
       GWSPH Research Day, April 15, 2026}
    };

  %% ============================================================
  %% ROW 1
  %% ============================================================

  \node[anchor=north west] at ($(TL)+(1.5cm,-14.8cm)$) {
    \begin{minipage}{56cm}
    \begin{posterblock}{Background \& Objectives}
      Shotgun metagenomics and amplicon sequencing frequently yield
      discordant community profiles, commonly attributed to biological
      or wet-lab differences --- but how much arises from
      \textsl{computational pipeline} choices?

      \medskip
      We isolate the dry-lab contribution by routing shotgun data
      through the amplicon pipeline, holding wet-lab inputs constant.
      We ask whether Fisher-Rao distance detects pipeline-driven
      differences more reliably than Bray-Curtis.

      \medskip
      \textbf{Dataset:} P\'{e}rez-Losada et al.\ (2022) vermicomposting;
      12 microcosms, $2\times2$ factorial (earthworm $\times$ time);
      3 replicates per cell. 554 ASVs; 99.1\% phylum assignment (SILVA 138.1).

      \medskip
      \begin{center}
      \includegraphics[width=0.95\linewidth]{../figures/pipeline_poster.pdf}
      \end{center}

      \smallskip
      {\small
      \textbf{\color[RGB]{70,110,150}Stream 1 --- Native Amplicon:}
      16S V4 PCR $\to$ Illumina $\to$ DADA2.\\[0.3ex]
      \textbf{\color[RGB]{140,90,156}Stream 2 --- MEGAHIT Contigs:}
      Shotgun $\to$ MEGAHIT $\to$ Barrnap $\to$ BBMap $\to$ DADA2.\\[0.3ex]
      \textbf{\color[RGB]{60,140,90}Stream 3 --- Read-Based:}
      Shotgun reads extracted and routed through amplicon DADA2.\\[0.3ex]
      Streams 2 \& 3 share identical wet-lab inputs, isolating
      computational pipeline as the variable of interest.}
    \end{posterblock}
    \end{minipage}
  };

  \node[anchor=north west] at ($(TL)+(59cm,-14.8cm)$) {
    \begin{minipage}{46.2cm}
    \begin{posterblock}{PCoA Ordination by Treatment $\times$ Time}
      \begin{center}
      \includegraphics[width=0.97\linewidth]{../figures/pcoa_permanova_phylum.pdf}
      \end{center}
      \begin{center}
      {\small Fisher-Rao (left) and Bray-Curtis (right) for
      Amplicon and MEGAHIT Contigs. Treatment and time structure
      visible under both; Fisher-Rao gives tighter clustering.}
      \end{center}
    \end{posterblock}
    \end{minipage}
  };

  %% ============================================================
  %% ROW 2
  %% ============================================================

  \node[anchor=north west] at ($(TL)+(1.5cm,-48.1cm)$) {
    \begin{minipage}{56cm}
    \begin{posterblock}{Phylum-Level Composition by Treatment $\times$ Time}
      \begin{center}
      \includegraphics[width=0.95\linewidth]{../figures/three_stream_experimental_design_barplot.pdf}
      \end{center}
      \begin{center}
      {\small Rows: analytical stream. Columns: experimental group
      (CT/EW $\times$ Day~21/63), 3 replicates per cell.
      Amplicon and MEGAHIT Contigs show concordant treatment response;
      Read-Based diverges ($<$0.1\% of shotgun reads match V4).
      FR distances below quantify the difference.}
      \end{center}
    \end{posterblock}
    \end{minipage}
  };

  \node[anchor=north west] at ($(TL)+(59cm,-53cm)$) {
    \begin{minipage}{46.18cm}
    \begin{posterblock}{Pairwise Distances: Fisher-Rao vs.\ Bray-Curtis}
      \begin{center}
      \includegraphics[width=0.97\linewidth]{../figures/pairwise_distances_boxplot.pdf}
      \end{center}
      \begin{center}
      {\small FR confirms concordance ($\bar{d}=0.693$) and divergence
      ($3\times$ larger for Read-Based), with lower variance than
      BC (CV ratio $\approx 0.44$) --- consistent with the composition above.}
      \end{center}
    \end{posterblock}
    \end{minipage}
  };

  %% ============================================================
  %% ROW 3
  %% ============================================================

  \node[anchor=north west] at ($(TL)+(1.5cm,-90.2cm)$) {
    \begin{minipage}{56cm}
    \begin{posterblock}{Fisher-Rao Distance \& PERMANOVA Results}

      \small
      Fisher-Rao distance arises necessarily from the Fisher
      Information Matrix and is the \textsl{unique} metric on the
      probability simplex invariant under sufficient statistics
      ({\v C}encov, 1982). Bray-Curtis was adopted by historical
      convention, not derived from the multinomial model.

      \medskip
      \begin{minipage}[t]{0.44\linewidth}
        For $\mathbf{p}$, $\mathbf{q}$ over $K$ taxa:
        \[
          D_{\mathrm{FR}} = 2\sqrt{K}\arccos\!\left(\sum_{k=1}^K \sqrt{p_k q_k}\right)
        \]
      \end{minipage}%
      \hfill
      \begin{minipage}[t]{0.52\linewidth}
        \vspace{0.3ex}
        \textbf{PERMANOVA} (9,999 permutations, phylum level):\\[0.5ex]
        \begin{tabular}{llrr}\hline
          & & \multicolumn{2}{c}{$p$-value} \\
          Stream & Effect & FR & BC \\ \hline
          Amplicon  & Treatment & 0.004** & 0.007** \\
                    & Time      & \textbf{0.073.} & 0.190 \\[0.3ex]
          MEGAHIT   & Treatment & 0.0001*** & 0.0004*** \\
                    & Time      & 0.079. & 0.093. \\ \hline
        \end{tabular}\\[0.5ex]
        \textit{FR detects time effect ($p=0.073$) missed by BC.}
      \end{minipage}

    \end{posterblock}
    \end{minipage}
  };

  \node[anchor=north west] at ($(TL)+(59cm,-82cm)$) {
    \begin{minipage}{46.18cm}
    \begin{posterblock}{Conclusions \& Future Directions}
      Much reported shotgun-amplicon discordance reflects
      \textsl{dry-lab pipeline differences}.\\%, not biology.
      Routing shotgun reads through the amplicon pipeline
      recovers concordant community profiles.

      \medskip
      Fisher-Rao distance, derived necessarily from the Fisher
      Information Matrix of the multinomial distribution, provides
      more stable community comparisons than Bray-Curtis.

      \medskip
      Further study of the sampling distribution of Fisher-Rao
      distance will support power analysis and Type~I error
      characterization for PERMANOVA-based testing.

      \medskip
      Validation against mock communities with known composition
      will establish the accuracy of contig-based 16S extraction
      across sequencing platforms.

      \medskip
      \begin{minipage}[t]{0.62\linewidth}
        \raggedright
        {\small
        \textbf{Data:} P\'{e}rez-Losada et al.\ (2022),
        \textit{Front.\ Microbiol.}\ 13:854423.
        BioProject PRJNA777435.\\[0.4ex]
        \textbf{Computing:} Pegasus HPC cluster,
        GW Research Technology Services.\\[0.4ex]
        \textbf{Theory:} \v{C}encov (1982), \textit{Statistical
        Decision Rules and Optimal Inference};
        Rao (1945).\\[0.4ex]
        \textbf{Manuscript:} Claude (Anthropic).}
      \end{minipage}%
      \hfill
      \begin{minipage}[t]{0.33\linewidth}
        \centering
        {\small\color{gray}\textbf{Poster \& Draft}}\\[0.3ex]
        
\includegraphics[width=4cm]{../figures/qrcode.pdf}\\[0.2ex]
        {\tiny github.com/cgaylord/\\2026GWSPH\_ResearchDayPoster}
      \end{minipage}
    \end{posterblock}
    \end{minipage}
  };

\end{tikzpicture}

\end{frame}
\end{document}